\section{Tablas de Mediciones}

\begin{table}[H]
\centering
\begin{tabularx}{\textwidth}{|X|Y|Y|Y|}
\hline
\textbf{Parámetro} & \textbf{Núcleo Macizo} & \textbf{Núcleo Laminado} & \textbf{Unidades} \\ \hline
Sección transversal del núcleo ($A$) & & & \\ \hline
Longitud media del camino magnético ($l$) & & & \\ \hline
Número de espiras devanado primario ($N_1$) & & & \\ \hline
Número de espiras devanado secundario ($N_2$) & & & \\ \hline
Resistencia del devanado ($R$) & & & \\ \hline
Factor de apilamiento & & & \\ \hline
\end{tabularx}
\caption{Parámetros geométricos y eléctricos de los núcleos ferromagnéticos.}
\label{tab:parametros}
\end{table}

\begin{table}[H]
\centering
\scriptsize
\begin{tabularx}{\textwidth}{| >{\hsize=1.4\hsize}Y | >{\hsize=0.955\hsize}Y | >{\hsize=0.955\hsize}Y | >{\hsize=0.955\hsize}Y | >{\hsize=0.955\hsize}Y | >{\hsize=0.955\hsize}Y | >{\hsize=0.955\hsize}Y | >{\hsize=0.955\hsize}Y | >{\hsize=0.955\hsize}Y | >{\hsize=0.955\hsize}Y |}
\hline
\textbf{Tipo de Núcleo} & \textbf{$V_{\text{RMS}}$ [V]} & \textbf{$I_{\text{RMS}}$ [A]} & \textbf{$P_{\text{medida}}$ [W]} & \textbf{$A_{\text{ciclo}}$} & \textbf{$K_h$} & \textbf{$K_v$} & \textbf{$P_{\text{núcleo}}$ [W]} & \textbf{$P_h$ [W]} & \textbf{$P_f$ [W]} \\ \hline
Macizo & & & & & & & & & \\ \hline
Laminado & & & & & & & & & \\ \hline
\end{tabularx}
\caption{Mediciones del circuito y parámetros del ciclo de histéresis.}
\label{tab:mediciones}
\end{table}

\begin{table}[H]
\centering
\small
\begin{tabularx}{\textwidth}{| c | c | Y | Y | Y | Y |}
\hline
\textbf{Tensión (V)} & \textbf{Corriente (A)} & \textbf{$B_{\text{max}}$ (Tesla)} & \textbf{$H_{\text{max}}$ (A/m)} & \textbf{Ancho del ciclo (T)} & \textbf{Área del ciclo (T $\cdot$ A/m)} \\ \hline
40 & & & & & \\ \hline
50 & & & & & \\ \hline
60 & & & & & \\ \hline
70 & & & & & \\ \hline
80 & & & & & \\ \hline
90 & & & & & \\ \hline
100 & & & & & \\ \hline
110 & & & & & \\ \hline
120 & & & & & \\ \hline
\end{tabularx}
\caption{Ciclo de histéresis para el Núcleo Macizo.}
\label{tab:histeresis_macizo}
\end{table}

\begin{table}[H]
\centering
\small
\begin{tabularx}{\textwidth}{| c | c | Y | Y | Y | Y |}
\hline
\textbf{Tensión (V)} & \textbf{Corriente (A)} & \textbf{$B_{\text{max}}$ (Tesla)} & \textbf{$H_{\text{max}}$ (A/m)} & \textbf{Ancho del ciclo (T)} & \textbf{Área del ciclo (T $\cdot$ A/m)} \\ \hline
40 & & & & & \\ \hline
50 & & & & & \\ \hline
60 & & & & & \\ \hline
70 & & & & & \\ \hline
80 & & & & & \\ \hline
90 & & & & & \\ \hline
100 & & & & & \\ \hline
110 & & & & & \\ \hline
120 & & & & & \\ \hline
\end{tabularx}
\caption{Ciclo de histéresis para el Núcleo Laminado.}
\label{tab:histeresis_laminado}
\end{table}
